\chapter{Fazit}

Nach der Beschreibung des Projektes sollte noch ein kurzes Fazit zum Projekt
selbst und der erstellten Applikation gegeben werden. So war es für dieses
Projekt möglich eine vollwertige Webapplikation für das die asynchrone Kommunikation
in Teams zu erstellen. So konnten alle wichtigen Komponenten wie die Nutzerverwaltung,
das Arbeiten mit den Teams und deren Kanälen und ein Echtzeit-Benachrichtigung-System
für die Textnachrichten und weitere Events implementiert werden. Insgesamt kann das
Projekt als Erfolg betrachtet werden.

Durch das Projekt konnte für mich perönlich gezeigt werdem, wie effektiv mit der
Programmiersprache Koltin Backends erstellt werden können. So wird im Vergleich
die Ergonomie zu der Programmiersprache Java verbessert, wobei sich diese auch in den
letzten Java Versionen verbessert hat. Auch können alle bisher geschriebenen Bibliotheken
für die JVM, wie \zb{} die für \emph{redis} oder \emph{Nats} einfach Verwendet und
integriert werden.

Eine weitere wichtige Entscheidung, ohne die ein hohes Maß an Interaktivität in
der Applikation schwierig gewesen wäre, war die Erstellung einer SPA. So musste
nicht das Templating der HTML-Dateien vom Backend übernommen werden. Dadurch
konnten einfach mit Hilfe von \emph{SolidJS} dynamische Anwendungen mit einem
hohen Maß an Interaktivität erstellt werden, wecher auch ihre Daten dank einer
Websocket-Verbindungn in Echtzeit anpassen können.

Zum Abschluss kann gesagt werden, dass mir die Entwicklung des Projektes eine
große Freude bereitet und einen tieferen Einblick in die Webentwicklung gegeben
hat.
